\DocumentMetadata{
  pdfversion=2.0,
  pdfstandard=ua-2,
  lang=en-US,
  tagging=on,
  tagging-setup={math/setup=mathml-SE,tabsorder=structure}
}
\documentclass[11pt,letterpaper]{article}
\usepackage[margin=0.68in,includefoot]{geometry}
\usepackage{newcomputermodern}
\usepackage{amsmath,amscd,mathtools}
\usepackage{tikz,adjustbox}
\providecommand{\square}{\mdlgwhtsquare}
\usepackage[hidelinks]{hyperref}
\hypersetup{
  pdftitle={Analysis First-Year Exam, January 2016, Part 1},
  pdfauthor={University of Florida Department of Mathematics},
  pdfsubject={Graduate mathematics examination},
  pdfdisplaydoctitle=true,
  bookmarksopen=true
}
\setcounter{secnumdepth}{0}
\setlength{\parindent}{0pt}
\setlength{\parskip}{0.28em}
\setlength{\emergencystretch}{3em}
\setlength{\leftmargini}{2.15em}
\setlength{\leftmarginii}{2.65em}
\setlength{\leftmarginiii}{3em}
\pagestyle{empty}
\raggedbottom
\allowdisplaybreaks
\NewTaggingSocketPlug{block/list/label}{examproblem}{%
  \tagstructbegin{tag=Lbl}\tagstructend%
  \tagstructbegin{tag=\LBody}%
  \tagstructbegin{tag=\ProblemHeadingTag,title-o={\ProblemHeadingText}}%
  \tagmcbegin{tag=\ProblemHeadingTag,actualtext={\ProblemHeadingText}}%
  #2%
  \tagmcend%
  \tagstructend%
}
\newenvironment{examproblems}[2]{%
  \def\ProblemHeadingTag{#2}%
  \AssignTaggingSocketPlug{block/list/label}{examproblem}%
  \begin{enumerate}%
  \setcounter{enumi}{#1}%
  \renewcommand{\labelenumi}{\textbf{\theenumi.}}%
  \setlength{\topsep}{0.35em}%
  \setlength{\partopsep}{0pt}%
  \setlength{\itemsep}{0.48em}%
  \setlength{\parsep}{0.18em}%
}{\end{enumerate}}
\newenvironment{examsubparts}{%
  \AssignTaggingSocketPlug{block/list/label}{default}%
  \begin{enumerate}%
  \setlength{\topsep}{0.18em}%
  \setlength{\partopsep}{0pt}%
  \setlength{\itemsep}{0.16em}%
  \setlength{\parsep}{0.1em}%
}{\end{enumerate}}
\newenvironment{examinstructions}{%
  \par\begingroup\itshape
}{%
  \par\endgroup\vspace{0.18em}
}
\makeatletter
\renewcommand\subsection{\@startsection{subsection}{2}{0pt}%
  {-1.25ex plus -.25ex minus -.1ex}%
  {0.55ex plus .1ex}%
  {\normalfont\large\bfseries}}
\renewcommand{\@maketitle}{%
  \newpage\null\vskip -1em%
  {\footnotesize\bfseries
    \href{https://www.ufl.edu/}{University of Florida}, \href{https://math.ufl.edu/}{Department of Mathematics}\par}%
  \vskip -0.5em%
  \tagstructbegin{tag=H1,title={Analysis First-Year Exam, January 2016, Part 1}}%
  \tagpdfparaOff%
  \tagmcbegin{tag=H1}%
    {\LARGE\bfseries\href{https://gma.math.ufl.edu/exams/analysis-first-year/}{Analysis First-Year Exam}, January 2016, Part 1\par}%
  \tagmcend%
  \tagpdfparaOn%
  \tagstructend%
  \vskip 0.25em%
  \tagmcbegin{artifact}%
    \hrule height 1pt%
  \tagmcend%
  \vskip 1em%
}
\makeatother
\title{Analysis First-Year Exam, January 2016, Part 1}
\author{}
\date{}
\begin{document}
\maketitle
\thispagestyle{empty}
\begin{examinstructions}
Do exactly four problems. Work should be presented in a logical fashion in order to receive full credit.
\end{examinstructions}
\begin{examproblems}{0}{H2}
\def\ProblemHeadingText{Problem 1.}
\item
Let \(A_1,\ldots,A_n,\ldots\) be subsets of a metric space~\(X\). If \(B=\bigcup_{i=1}^n A_i\), prove that \(\overline{B}=\bigcup_{i=1}^n \overline{A_i}\), where \(\overline{B}\) stands for the closure of~\(B\). If \(C=\bigcup_{n=1}^{\infty}A_n\), prove that \(\bigcup_{i=1}^{\infty}\overline{A_i}\subset\overline{C}\).
\def\ProblemHeadingText{Problem 2.}
\item
Let~\(f\) be a mapping from \(B_1\) to itself, where \(B_1\) is the unit ball in \(\mathbb{R}^n\) (\(n\geq 1\)). Let \(R(f)\) be the range of~\(f\) in \(B_1\). Prove that~\(f\) is continuous if and only if for any open set~\(V\) in \(B_1\), \(f^{-1}(V\cap R(f))\) is also open. Your reasoning should be based on the following definition of continuity: \(f\) is called to be continuous at~\(x\) if for any given \(\varepsilon>0\), there exists \(\delta>0\) such that all \(y\in B(x,\delta)\) (the ball centered at~\(x\) with radius~\(\delta\)) satisfies \(|f(y)-f(x)|<\varepsilon\).
\def\ProblemHeadingText{Problem 3.}
\item
Let~\(K\) be a compact subset of \(\mathbb{R}^n\) and~\(F\) be a closed subset of \(\mathbb{R}^n\) (\(n\geq 1\)). Suppose \(K\cap F=\varnothing\), prove that
\[
\operatorname{distance}(K,F):=\inf_{x\in K,\,y\in F}|x-y|>0.
\]
 Moreover, there exist \(x_0\in K\) and \(y_0\in F\) such that \(|x_0-y_0|=\operatorname{distance}(K,F)\).
\def\ProblemHeadingText{Problem 4.}
\item
Let~\(f\) be a monotone function on \((0,1)\). Prove that the number of points where~\(f\) is discontinuous is at most countable.
\def\ProblemHeadingText{Problem 5.}
\item
Let~\(f\) be continuous function from \([0,2]\) to \([0,2]\). Prove that~\(f\) has a fixed point in \([0,2]\). Give an example to show that if~\(f\) is not continuous, the conclusion may not hold.
\def\ProblemHeadingText{Problem 6.}
\item
Suppose \(\sum_{n=1}^{\infty}a_n\) is convergent and every \(a_n\geq 0\) (\(n=1,\ldots,\infty\)), prove that \(\sum_{n=1}^{\infty}\frac{\sqrt{a_n}}{n^{2/3}}\) is convergent. Give an example to show that \(\sum_{n=1}^{\infty}\frac{a_n^{1/4}}{n^{2/3}}\) could be divergent.
\end{examproblems}
\end{document}
